\subsubsection{Coordonnées centrales, orientation et registres régionaux}
\label{el-centro:registros}

La structure centrale relie position, évaluation et transport. Dans la
cellule \((5,5)\), les deux feuillets d'APP donnent
\begin{equation}
 \mathcal A(5,5)=\bigl((1,1),(7,2)\bigr),\qquad
 5+5=1+9\cdot1,\qquad5\cdot5=7+9\cdot2.
 \label{el-centro:app}
\end{equation}
Les deux unités du feuillet additif sont, dans cet ordre, le résidu et le quotient.
Leur coïncidence numérique conserve deux fonctions arithmétiques distinctes.
Sur la fibre de somme \(10\), les deux lacunes orientées vérifient
\begin{equation}
 \mathcal A(8,2)=\mathcal A(2,8)=\bigl((1,1),(7,1)\bigr).
 \label{el-centro:vacancias}
\end{equation}
Le feuillet additif complet et le résidu multiplicatif sont conservés ;
la coordonnée de quotient du produit diminue exactement de
\(1\). Le transport retient également laquelle des deux positions
conjuguées a été atteinte.

\begin{figure}[H]
\centering
\begin{tikzpicture}[x=.67cm,y=.67cm,font=\footnotesize,>=Latex]
 \draw[step=1,black!12,line width=.3pt] (1,1) grid (9,9);
 \draw[->,black!60] (.7,1)--(9.5,1) node[right] {$i$};
 \draw[->,black!60] (1,.7)--(1,9.5) node[above] {$j$};
 \foreach \k in {1,...,9}{
  \node[below=2pt] at (\k,1) {$\k$};
  \node[left=2pt] at (1,\k) {$\k$};
 }
 \draw[black!65] (1,9)--(9,1);
 \draw[<->,line width=.9pt] (2,8)--(8,2);
 \foreach \x/\y in {2/8,5/5,8/2}{
  \filldraw[fill=white,line width=.8pt] (\x,\y) circle (3pt);
 }
 \fill (5,5) circle (1.5pt);
 \node[above right=3pt,fill=white,inner sep=1pt] at (2,8) {$(2,8)$};
 \node[above right=3pt,fill=white,inner sep=1pt] at (5,5) {$(5,5)$};
 \node[above right=3pt,fill=white,inner sep=1pt] at (8,2) {$(8,2)$};
 \node[anchor=west,align=left] at (10.5,7.6)
   {$i+j=10$\\$R^+=1,\quad q^+=1$};
 \node[anchor=west,align=left] at (10.5,5.4)
   {centre : $ij=25$\\$R^\times=7,\quad q^\times=2$};
 \node[anchor=west,align=left] at (10.5,3.1)
   {lacunes : $ij=16$\\$R^\times=7,\quad q^\times=1$};
\end{tikzpicture}
\caption{Centre électronique et vacances conjuguées dans la coordinatisation APP.
La diagonale est la fibre additive \(\mathcal F_{10}\) ; le centre
est le point fixe de \((i,j)\mapsto(10-i,10-j)\). Les deux
déplacements de module \(3\) conservent la somme et le reste
du produit et réduisent d'une unité le quotient multiplicatif. Les
flèches représentent les déformations sur la fibre.}
\label{el-centro:figura}
\end{figure}


La symétrie centrale caractérise aussi sa persistance lors des changements de
coordinatisation. Si une transformation admissible \(f\) de l'atlas entrelace
l'inversion, \(f\rho=\rho f\), alors
\[
 \rho f(5,5)=f\rho(5,5)=f(5,5).
\]
L’unicité du point fixe impose \(f(5,5)=(5,5)\). C’est la
propriété interne qui rend stable la localisation centrale : toute
transformation de ce type conserve la cellule avant d’évaluer son
caractère électronique (section~\ref{el-centro:registros}).

Le bloc matriciel central appartient à l'évaluation sur les deux
positions adjacentes \(4,5\). Sa décomposition, établie dans le
développement local du noyau, est
\[
 B_c=\begin{pmatrix}7&2\\2&7\end{pmatrix}=9P_++5P_-.
\]
Ici, \(9\) et \(5\) sont les valeurs propres de deux modes, tandis que
\((5,5)\) est une position de l'atlas. Les concaténations de ses
entrées donnent \(72+27=77+22=99\) et
\(77-22=55=54+1\). La première égalité conserve la somme totale sous
deux dispositions ; la seconde exprime le défaut entre les lectures
diagonale et antidiagonale. Ces opérations précèdent la composition
électronique et conservent leur provenance matricielle.

L'orientation initiale est une autre donnée. Dans la représentation du
retour électronique, on fixe \((h_0,v_0)=(1,1)\), c'est-à-dire une avance
positive dans les deux coordonnées. L'inversion simultanée tous les \(27\)
pas, lue par fenêtres de \(9\), détermine
\[
 \tau_e=(+1,+1,+1,-1,-1,-1,+1,+1,+1,-1,-1,-1).
\]
Les segments \((1,1,1)\) appartiennent à cette suite d'orientations :
chaque segment comprend trois fenêtres de même sens. La récurrence
de la section consacrée aux opérateurs de lecture démontre ce registre dans
\eqref{eq:el-tau}, et ses deux évaluations restituent le triplet
\((80,54,6)\) qui intervient dans
\eqref{eq:el-formula-completa}. La coordonnée de position, les
quotients APP et la suite des signes sont ainsi reliés par
des opérations explicites.

La signature \(555555\) appartient à une lecture régionale d'un autre
domaine. Pour une émission décimale \(U=(U_1,\ldots,U_6)\), on écrit
\(U_j=100d_{1j}+10d_{2j}+d_{3j}\), avec trois chiffres, y compris les
zéros initiaux. L'opérateur de lecture de la signature multiplicative est
\[
 \mathcal E_\times(U)
   =\bigl([d_{2j}-d_{1j}]_{10}\bigr)_{j=1}^{6}.
\]
Dans la région transversale
\(U_\pi^\perp=(501,614,498,169,272,272)\), les différences sont
\((-5,-5,5,5,5,5)\). Par conséquent,
\begin{equation}
 \mathcal E_\times(U_\pi^\perp)=(5,5,5,5,5,5),\qquad
 \Psi(U_\pi^\perp)=010211.
 \label{el-centro:firma-transversal}
\end{equation}
La première application conserve une signature multiplicative du catalogue ;
la seconde conserve le mot ternaire commun aux cinq régions.
La figure \ref{pi-estrella:figura} montre la position transversale de
cette émission et sa multiplicité \(432\).

\paragraph{Mémoire nécessaire au transport de la signature constante.}
La signature \(555555\) fournit également un exemple vérifiable de
l'information de trajectoire qu'exige TPK. Soit
\(\mathcal X_\times=D_9^2\times\{N,E,S,O\}\), avec \(324\)
positions orientées du canal produit. L'avance \(P\) déplace
d'une case dans la direction mémorisée et le demi-tour \(H\)
inverse cette direction. Dans cette représentation, l'opérateur de lecture
\(E_{\rm sig}\) somme par fenêtres les exposants discrets
\[
 \ell_9(1,2,4,8,7,5)=(0,1,2,3,4,5),\qquad
 \ell_9(3)=\ell_9(6)=\ell_9(9)=0.
\]
Pendant chacune des six fenêtres de neuf pas, on évalue
\(\ell_9(\rho_9(ij))\) avant chaque avance produit, sélectionnée
par TRIT aux instants \(2,5,8\pmod9\). La somme de la fenêtre est
réduite modulo \(10\). Après les pas \(27\) et \(54\), on inverse
la direction. Cette règle détermine les six composantes de
\(E_{\rm sig}\) à partir de la position et de l'orientation initiales.

L'énumération de ce domaine produit \(26\) signatures. La fibre de
\(555555\) contient \(24\) positions orientées ; après une
avance, leurs images se répartissent en trois signatures, avec huit
représentants chacune :
\[
 555177,\qquad555717,\qquad555771.
\]
Trois témoins, dans les coordonnées de \(1\) à \(9\), permettent de vérifier
directement la distinction :
\[
\begin{array}{c|c|c}
x&E_{\rm sig}(x)&E_{\rm sig}(Px)\\\hline
(3,5,N)&555555&555177\\
(3,4,N)&555555&555717\\
(3,4,S)&555555&555771
\end{array}
\]
La valeur commune de la signature ne détermine donc pas sa valeur suivante :
la classe de transport à l'intérieur de cette fibre fait partie de la donnée
opératoire. La partition stable minimale se calcule à partir de
l'équivalence des signatures \(\sim_0\), en raffinant par
\[
 x\sim_{n+1}y\quad\Longleftrightarrow\quad
 x\sim_n y,\quad Px\sim_n Py,\quad Hx\sim_n Hy.
\]
Le procédé s’achève parce que le domaine est fini. Toute congruence
stable qui raffine \(\sim_0\) raffine chaque \(\sim_n\), par récurrence ;
son point fixe est donc la congruence stable minimale requise.
Le recensement donne \(26\to28\to28\), avec l’histogramme
\(27\cdot8+108=324\). La seule signature qui se décompose est
\(555555\), en trois classes de huit éléments. Le certificat joint
reproduit l’énumération, les trois classes et la stabilité par les deux
opérateurs (section~\ref{el-centro:registros}).

La multiplicité \(24\) appartient ici à la fibre du canal produit ;
la multiplicité \(432\) appartient à l'émission régionale complète
\(U_\pi^\perp\). La première démontre pourquoi la signature exige une mémoire
de transport ; la seconde situe cette signature dans le catalogue conjoint
des deux feuillets. Les deux lectures participent de la même architecture,
avec des domaines et des opérations spécifiés.

Cette distinction des domaines permet de réunir les occurrences de l'unité
et du centre avec leurs significations respectives :
\begin{center}
\small
\begin{tabularx}{\textwidth}{@{}lX@{}}
\toprule
Registro & Objet et fonction \\
\midrule
\((5,5)\) & Cellule fixe de l'inversion centrale de l'atlas APP.\\
\((R^+,q^+)=(1,1)\) & Résidu et quotient de l'évaluation additive centrale.\\
\((h_0,v_0)=(1,1)\) & Orientation initiale des deux déplacements.\\
\((+1,+1,+1)\) & Trois fenêtres consécutives de \(\tau_e\) d'orientation positive.\\
\(c=111111\) & Donnée de la cinquième frontière nonadique, section \ref{mono-recta:desarrollo}.\\
\(\mathcal E_\times=555555\) & Signature de la région transversale de \(\pi\), à six positions.\\
\bottomrule
\end{tabularx}
\end{center}

La coordination électronique s'obtient en composant ces niveaux dans leur
ordre de dépendance. APP détermine le centre et les évaluations ;
TRIT sélectionne régimes et orientations ; TPK prolonge la trajectoire
et conserve ses registres. Les régions préalablement engendrées apportent
les caractères \(\pi,\varphi,e\) ; le registre dodécaphasique apporte
\(\alpha\) ; les opérateurs de lecture électroniques déterminent les facteurs de
l'équation \eqref{eq:el-formula-completa}. Cette composition conserve
l'identité de la section à travers ses différentes
représentations, jusqu'à l'évaluation dimensionnelle et à la comparaison
métrologique ultérieure.
